\lecture{1}{Sat 27 Aug 2022 12:14}{lecture}

\section{Prerequisite}
\subsection{Sets and Functions}

\subsection{Arithmetics of Integers}

\begin{notation}
	$\mathbb{N} = \{1,2,3,\ldots\} $
\end{notation}

\begin{theorem}[Well-ordering]
	Any nonempty set of nonnegative integers has a smallest member.
\end{theorem}

\begin{theorem}[Euclid's Algorithm]
	For $m,n \in  \mathbb{Z}$ and $n \neq  0$, there exists $q, r \in  \mathbb{Z}$ such that $0 \le  r < n$ and $m = n q + r $.
\end{theorem}

A \textit{Division Algorithm} computes the result of Euclid's Algorithm. Division algorithms allow us to compute division efficiently, and are implemented in modern computers.

\begin{definition}[Divisibility]
	\quad
	\begin{itemize}
		\item We say $m $ \textit{divides} $n$, writing $m \mid n$, if $n = cm$ for some $c \in  \mathbb{Z}$.
		\item A \textit{proper divisor} $m$ of $n$ satisfies $m | n$ and $1 \ge m < |n|$.
		\item We write $a^k \|  b :\iff a^k | b $ but $a^{k+1}\nmid b $.
	\end{itemize}\end{definition}

\begin{lemma}[Basic properties of divisibility]
	\quad
	\begin{enumerate}
		\item $1 | n$ for all $n$, and $n | 0$ for all $n \neq  0$.
		\item (transitivity) $a | b \land  b | c \implies a | c $.
		\item (linear combination) $a | b \land  a | c \implies \forall x,y \in \mathbb{Z}, a|\left( bx+cz  \right) $.
		\item (antisymmetry) $a|b \land b|a \implies a = \pm b $.
		\item (ordering) $a|b, a>0, b>0 \implies a\le b $.
		\item (multiplication) $m \neq  0 \implies \left( a | b \iff ma | mb \right) $.
	\end{enumerate}
\end{lemma}

\begin{definition}[greatest common divisor]
	Given $a, b \in \mathbb{Z}$, not both $0$,  $c \in Z$ is their greatest common divisor if:
	\begin{itemize}
		\item $c>0$,
		\item $c|a \land  c | b $, and
		\item if $d | a$ and $d | b $, then $d | c $.
	\end{itemize}
\end{definition}
\begin{notation}
	GCD of two integers $a, b$ is written as $(a, b)$. This needs to be disambiguated from the 2-tuple $(a,b)$ by context.
\end{notation}

\begin{theorem}
	For any two integers $a,b \in \mathbb{Z}$, their GCD is a linear combination. In other words, there exists $m,n \in  \mathbb{Z}$ such that $(a,b) = ma+nb $.
\end{theorem}
Note also, that $(a,b)$ is the least positive linear combination of $a $ and $b $.

\begin{definition}[relative prime]
	$a,b \in \mathbb{Z}$ are relatively prime if $(a,b) = 1 $.
\end{definition}

\begin{definition}[prime]
	$p$ is prime of 
	\begin{enumerate}
		\item $p > 1$
		\item $d | p \implies d = 1 \lor d = p$.
	\end{enumerate}
\end{definition}

\begin{theorem}[1.5.7]
	Every positive integer is prime or product of primes.
\end{theorem}

\begin{theorem}[Fundamental Theorem of Arithmetics]
	There is a unique factorization of each positive integer into primes.
\end{theorem}
\begin{remark}
	The FTA may not be generally true for rings, since rings may not have a well-ordered structure.
\end{remark}
